\documentclass[11pt,a4paper]{article}

\usepackage[T1,T2A]{fontenc}
\usepackage[utf8]{inputenc}
\usepackage[english,russian]{babel}
\usepackage{geometry}
\usepackage{hyperref}
\usepackage{enumitem}
\usepackage{csquotes}
\usepackage{booktabs}
\usepackage{amsmath}

\geometry{margin=2.2cm}
\setlist{nosep,leftmargin=1.4em}
\setlist[description]{leftmargin=0pt,itemsep=0.5em,style=unboxed}
\hypersetup{colorlinks=true,linkcolor=black,urlcolor=blue}
\setlength{\parskip}{0.45em}
\setlength{\emergencystretch}{3em}

\begin{document}

\begin{center}
{\LARGE\bfseries TenderAI: где мы сейчас и куда идём дальше}\\[0.7em]
{\large Постановка задачи, гипотезы, что уже закрыто в коде и что осталось}\\[0.5em]
{\small 6 ноября 2026}
\end{center}
\bigskip

\section{Что мы решаем}

На входе десять закупок и четыре артикула эластомерных инфузионных помп: \texttt{CT0050L}, \texttt{MC1248L}, \texttt{MS5248M} и \texttt{M275-07}. Нужно понять, какой артикул технически подходит под какую позицию закупки, и доказать это цитатами с обеих сторон. Цена, сроки, адрес поставки и документы участника в задачу не входят, только техника и комплектация.

Каждый артикул~--- это Excel с листом «Характеристики», всего 144 строки. Закупки устроены сложнее: 72 файла в форматах DOC, DOCX, XLS, XLSX, PDF, XML, TXT и RAR. Расширение иногда не совпадает с содержимым, три файла вместо документа содержат сохранённый ответ об ошибке ЕИС, а часть требований лежит не в таблицах, а в обычной прозе.

Ответ нужен на трёх уровнях: требование против артикула (900 пар), позиция против артикула (76 пар) и итоговая сводка «10 закупок на 4 артикула» (40 пар). На каждом уровне статус один из трёх.

\begin{description}
  \item[\texttt{match}.] Все применимые технические условия подтверждены, и у каждого есть источник.
  \item[\texttt{mismatch}.] Доказано нарушение обязательного условия.
  \item[\texttt{unknown}.] Всё остальное: значения нет, единицы несовместимы, требование не разобрано или в закупке битый файл ЕИС (\texttt{coverage\_gap}).
\end{description}

У задания есть несколько жёстких правил, и почти все наши решения из них вытекают. Близость по смыслу не равна соответствию: числа, диапазоны, множества и условия «и/или» проверяет код. Характеристики соседней модели на артикул переносить нельзя. Пара, которую поиск отсёк, не становится \texttt{mismatch} автоматически. Отсутствие сведений о болюсе или УФ-защите не значит, что свойства нет.

В рабочем пайплайне участвует \texttt{gpt-5.6-luna} (через \texttt{cursor-sdk}), а расход токенов и вызовов пишется по этапам в \texttt{usage.jsonl}. Главный акцент оценки~--- латентные представления: как артикул и требование переводятся в общее векторное пространство и насколько это помогает сопоставлению по сравнению с простым поиском по словам.

\section{Как мы проверяем идеи}

Чтобы гипотеза что-то решала, нужен общий измеритель. Мы вручную проверили и закрыли отложенную выборку в \nolinkurl{eval/held_out_draft.json}: 15 пар «требование на артикул» (\texttt{G01}--\texttt{G15}) по первичным файлам четырёх закупок и контрольный список \texttt{E01} из 12 прозовых требований первой позиции закупки 02. Каждую гипотезу проверяем одинаково: фиксируем вход, меняем один компонент и сравниваем метрики.

Метрики те же, что в задании: полнота и точность извлечения требований, \texttt{Recall@k} при поиске характеристики внутри SKU, \texttt{macro-F1} по трём статусам, точность класса \texttt{match}, доказательность цитат и токены по этапам.

\section{Гипотезы и их текущий статус}

Мы разбили их на четыре группы. Семь гипотез из десяти полностью закрыты и проверены в коде, а три онлайн-эксперимента ждут только свежего ключа \texttt{CURSOR\_API\_KEY}.

\subsection*{Представление и латентное пространство}

\begin{description}
  \item[H1. Кодировать нужно характеристику, а не изделие целиком.] Вектор на каждую характеристику против одного «мешка» на весь артикул. \emph{Подтвердилась: убирает 533 ложных срабатывания.}
  \item[H2. Латентное пространство FRIDA + RRF находит характеристику лучше лексики.] Плотные эмбеддинги \texttt{ai-forever/FRIDA} (1536d) с префиксами \texttt{search\_query:} / \texttt{search\_document:} и слиянием рангов через $\mathrm{RRF}(k=60)$ против лексического поиска. \emph{Подтвердилась: точность выравнивания выросла с 0,667 до 0,933, \texttt{Recall@3} достиг 1,000, ложные \texttt{match} упали до 0, а на полном корпусе открылись 2 полных \texttt{match} закупок.}
  \item[H3. Гибрид сохраняет точность там, где одна близость ошибается.] Статус по близости против схемы «поиск плюс проверка кодом». \emph{Подтвердилась: отсекает 156 близких, но несовместимых пар.}
\end{description}

\subsection*{Извлечение требований из грязных документов}

\begin{description}
  \item[H4. GPT-Luna достаёт требования из прозы, которую не берёт табличный парсер.] \emph{Подтвердилась: +118 требований, охват 10/10 закупок, на срезе \texttt{E01} точность и полнота 12/12.}
  \item[H5. Поиск фрагментов до вызова модели экономит токены без потери требований.] Весь документ против top-$k$ релевантных страниц на закупках 02 и 03 (где полный документ 03 стоил 42\,832 input-токенов). \emph{Ждёт живой \texttt{CURSOR\_API\_KEY}.}
  \item[H9. OCR нужен точечно, только там, где у PDF нет текстового слоя.] \emph{Подтвердилась и подключена в коде через Apple Vision ровно для 1 пустого PDF из 8.}
\end{description}

\subsection*{Кто выносит вердикт}

\begin{description}
  \item[H7. Код не хуже модели на числах и при этом дешевле.] Кодовый чекер против вердикта \texttt{gpt-5.6-luna} на 15 проверенных парах. \emph{Ждёт живой \texttt{CURSOR\_API\_KEY}.}
  \item[H8. Каскад держит качество за долю токенов.] Сначала код за 0 токенов, вызов \texttt{gpt-5.6-luna} только на \texttt{unknown} (например, составное требование \texttt{G14}). \emph{Ждёт живой \texttt{CURSOR\_API\_KEY}.}
  \item[H10. Нормализация единиц и обработка отрицаний несут нагрузку.] \emph{Подтвердилась: от отрицаний зависят 84 вердикта из 391, единицы страхуют 9 пар.}
\end{description}

\subsection*{То, что мы отклонили}

\begin{description}
  \item[H6. Реранкер поверх четырёх артикулов не окупается.] \emph{Подтвердилась: мы проверяем все 4 SKU ($\text{Recall@4}=1{,}0$), реранкер отклонён.}
\end{description}

\section{Что сделано и что показала проверка}

\subsection{Закрыты все три технических пробела в парсинге и два бага чисел}

\begin{itemize}
  \item \textbf{Закупка 09 (\texttt{02\_Извещение.xml}):} файл оказался настоящим XLSX с расширением \texttt{.xml}. Передача \texttt{io.BytesIO} в \texttt{openpyxl} сняла отказ по расширению.
  \item \textbf{Закупка 05 (\texttt{02\_5.\ пк.rar}):} подключена рекурсивная распаковка через \texttt{bsdtar}. Распаковка показала важный факт: внутри архива лежали только бланк проекта контракта и пустая таблица спецификации («не менее 6 месяцев» остаточного срока годности), а само внешнее ТЗ лежало в битом файле \texttt{08\_Электронный документ...docx} (\texttt{saved\_error\_response}). Значит, статус \texttt{unknown} (\texttt{coverage\_gap}) по закупке 05 абсолютно честный.
  \item \textbf{Закупка 04 (\texttt{06\_Другое..}):} подключён гейтированный OCR на Apple Vision (\texttt{ocr.py}, кэш по SHA-256). Распознавание показало, что внутри одностраничный скан декларации по ПП РФ №~1875 об отсутствии российского производства, то есть ни одна техническая характеристика помпы в закупке 04 не потеряна.
  \item \textbf{Исправления в \texttt{parse.py}:} одиночные верхние неравенства (\texttt{<=}, $\le$, \texttt{<}) теперь разбираются как верхняя граница \texttt{lte/lt} (раньше \texttt{$\le 2$ мл} остаточного объёма ошибочно превращалось в \texttt{gt 2.0}), а записи вида \texttt{2:4:6:8:10:12:14} разбираются как числовое множество \texttt{set}.
\end{itemize}

В итоге в \texttt{runs/full/errors.json} осталось ровно 3 записи~--- три реальных JSON-ответа об ошибке выгрузки ЕИС в закупках 02, 05 и 06.

\subsection{Ручная проверка отложенной выборки}

Вычитка \texttt{eval/held\_out\_draft.json} глазами показала, почему нельзя доверять совпадению общих сумм статусов. В лексическом baseline распределение было таким же, как в эталоне (\texttt{7 match / 4 mismatch / 4 unknown}), но в \texttt{G11} скрывался ложный \texttt{match} (номинальный объём заполнения притянулся к максимальному 300 мл у 250-миллилитровой помпы \texttt{CT0050L}), в \texttt{G12}~--- потерянный \texttt{match} («Регулирование скорости: нет» не сматчилось с «Тип скорости: Постоянная»), а общая точность выравнивания составляла всего 10 из 15 (\texttt{0,667}). Подробный пообъектный разбор оформлен в \nolinkurl{eval/REPORT.md} и в итоговом отчёте \nolinkurl{docs/final_report.tex}.

\subsection{Латентное пространство FRIDA и гибрид RRF (H2)}

При подключении \texttt{ai-forever/FRIDA} (1536d) выяснилась принципиальная деталь: если класть в текст энкодера сырые числа и флаги «Да/Нет», в косинусном пространстве возникают межсвойственные коллизии~--- например, \emph{«Скорость инфузии, мл/час: 5»} тянется к \emph{«Объём болюса (мл): 5»} ($\cos = 0{,}391$) сильнее, чем к \emph{«Настройки скорости: 2/4/6/8/10/12/14»} ($\cos = 0{,}174$).

Поэтому в \texttt{embeddings.py} мы кодируем предметно-оформленный \textbf{семантический слот} (\texttt{search\_query: / search\_document: Инфузионная помпа --- <характеристика[, ед.\ изм.][: категориальный текст]>}), а все числа, интервалы, множества и отрицания оставляем строго в \texttt{Constraint}. Ранг по косинусу FRIDA ($\ge 0{,}625$) объединяется с лексическим рангом через $\mathrm{RRF}(k=60)$ с фильтром метрологических квалификаторов (\emph{минимальный / номинальный / максимальный / остаточный}).

\begin{center}
\begin{tabular}{lrrr}
\toprule
Метрика на \texttt{held\_out\_draft.json} (15 пар) & Лексика & FRIDA dense & Гибрид RRF \\
\midrule
Точность выравнивания (\texttt{alignment\_accuracy}) & 0,6667 & 0,8667 & \textbf{0,9333} \\
\texttt{Recall@1} / \texttt{Recall@3} (12 пар со слотом) & 0,8333 / 0,8333 & 0,9167 / 1,0000 & \textbf{0,9167 / 1,0000} \\
Отказ на отсутствующих слотах (\texttt{G08, G11, G13}) & 0 / 3 & 2 / 3 & \textbf{3 / 3} \\
\texttt{macro-F1} по статусам & 0,8690 & 0,9373 & \textbf{0,9373} \\
Доказательная точность \texttt{match} & 0,7143 & 1,0000 & \textbf{1,0000} \\
Ложные \texttt{match} & 1 (\texttt{G11}) & 0 & \textbf{0} \\
\bottomrule
\end{tabular}
\end{center}

\subsection{Единый прогон по 10 закупкам и два подтверждённых \texttt{match}}

В папке \texttt{runs/full} собран полный прогон с заполненными 1536-мерными векторами FRIDA в \texttt{representations/structured.json} (369 объектов). На 900 парах доля \texttt{not\_aligned} снизилась с 45,2\% (407) до \textbf{35,6\% (320)}, а число разрешимых пар выросло с 391 до \textbf{445}.

Главное изменение на уровне позиций и закупок: если раньше подтверждённых совпадений было ноль, то теперь артикул \textbf{\texttt{CT0050L}} получил стопроцентно доказанный \textbf{\texttt{match} (7 из 7 требований)} сразу в двух закупках~--- \textbf{07 (\texttt{0372100048826000117})} и \textbf{10 (\texttt{0308300014226000025})}. В обеих закупках латентное пространство дотянуло две недостающие характеристики: \emph{«Путь введения: Внутривенный»} $\to$ \emph{«Внутривенное введение: Да»} ($\cos = 0{,}7358$) и \emph{«Регулирование скорости инфузии: Нет»} $\to$ \emph{«Тип скорости: Постоянная»} ($\cos = 0{,}6360$). Код покрыт 37 автотестами, все проходят.

\section{Что осталось доделать}

Все оффлайн-артефакты сдачи (архитектура, таблица обоснования компонентов, таблица «baseline против итога», матрица $10 \times 4$, векторы в \texttt{runs/full} и сценарий демонстрации) уже собраны в \nolinkurl{docs/final_report.tex}. При появлении активного ключа \texttt{CURSOR\_API\_KEY} остаётся прогнать три онлайн-замера с живыми вызовами \texttt{gpt-5.6-luna}:
\begin{enumerate}
  \item \textbf{H5:} подать в \texttt{gpt-5.6-luna} только релевантные страницы закупок 02 и 03 и сравнить входные токены и полноту с полным прогоном (42\,832 input-токенов на закупке 03).
  \item \textbf{H7:} сравнить вердикт \texttt{gpt-5.6-luna} с кодовым чекером на 15 проверенных парах по \texttt{macro-F1} и расходу токенов.
  \item \textbf{H8:} прогнать каскад «код за 0 токенов + \texttt{gpt-5.6-luna} только на \texttt{unknown}» (проверить, дорешивает ли каскад двухстрочный случай \texttt{G14} без потери точности).
\end{enumerate}

\subsection*{Как воспроизвести текущее состояние}

\begin{itemize}
  \item Единый прогон (FRIDA + RRF + OCR + RAR + кэш Luna):\\ \texttt{PYTHONPATH=src python -m tenderai run --dataset dataset\_zakupki --out runs/full}
  \item Абляции H1, H2, H3, H6, H9 и H10:\\ \texttt{python scripts/experiments.py}
  \item Автотесты (37 шт.): \texttt{pytest -v}
  \item Сводная таблица токенов: \texttt{python scripts/token\_table.py}
\end{itemize}

\end{document}
